% P17-D3: finite information value and Blackwell comparison.
\chapter{Value of Information and Comparison of Experiments}
\label{sec:p17-value-of-information}

Information has decision value when it permits a better choice under a
specified prior and payoff function. An experiment can have value for one
decision and none for another. A comparison that holds for every decision
problem is therefore stronger than a favorable calculation in one example.
This chapter gives both calculations and the full finite comparison theorem.

The source presentation is Blackwell's \emph{Equivalent Comparisons of
Experiments}, especially its finite matrix development and
Sherman--Stein theorem
\cite[pp.~265--268]{Blackwell1953EquivalentExperiments}.
The proof below supplies an explicit finite-dimensional separation argument
and also the source's trace/minimax route, with the required finite minimax
result proved. The posterior and finite-vector formulations are proved as
well. The general observation-space extension in the source has additional
measure-theoretic requirements.

\section{Experiments and decisions}
\label{sec:p17-experiment-decisions}

\begin{hbdefinition}{def:p17-finite-experiment}{States, signals and a finite experiment}
Let \(S=\{1,\ldots,n\}\) be a nonempty finite set of states and
\(Z=\{1,\ldots,r\}\) a nonempty finite set of signals.
A \emph{finite experiment} is a matrix \(P=(p_{sz})\) with
\[
 p_{sz}\geq0,\qquad \sum_{z=1}^r p_{sz}=1\quad(s\in S).
\]
Such a matrix is called \emph{row stochastic}, or a Markov matrix in
the source. Its row \(s\) specifies the conditional distribution of the
observed signal when the state is \(s\). It is a supplied model, rather
than a conclusion drawn from a sample.

A prior is a vector \(\pi=(\pi_s)_{s\in S}\) with nonnegative entries
summing to one. Its support is \(S_\pi=\{s:\pi_s>0\}\).
The joint and marginal probabilities are
\[
 \Pr(s,z)=\pi_s p_{sz},\qquad q_z=\sum_s\pi_s p_{sz}.
\]
For \(q_z>0\), define the posterior
\[
                  \pi^z_s=\frac{\pi_s p_{sz}}{q_z}.
\]
No posterior is assigned by this formula when \(q_z=0\).
The set of probability vectors on \(S\) is denoted \(\Delta(S)\).
\index{experiment!finite information structure}
\index{posterior probability!finite experiment}
\end{hbdefinition}

\begin{hbdefinition}{def:p17-information-policy}{Actions, policies and expected value}
Let \(A\) be a nonempty finite action set and let \(u(a,s)\in\mathbb R\)
be the payoff of action \(a\) in state \(s\).
A randomized policy \(\delta\) has entries \(\delta(a\mid z)\geq0\)
with \(\sum_a\delta(a\mid z)=1\) at every signal.
Its randomization is conditional on the signal, with no additional state
information. A deterministic policy selects one action at each signal.
Define
\[
 \begin{aligned}
 J(\pi,P,u,\delta)
   &=\sum_{s,z,a}\pi_s p_{sz}\delta(a\mid z)u(a,s),\\
 V(\pi,P,u)&=\max_\delta J(\pi,P,u,\delta),\\
 V_0(\pi,u)&=\max_{a\in A}\sum_s\pi_s u(a,s),\\
 V_{\mathrm{PI}}(\pi,u)&=\sum_s\pi_s\max_{a\in A}u(a,s).
 \end{aligned}
\]
The next theorem proves the existence of the first maximum.
The last expression is the value with the state itself observed, called
perfect information.
Payoffs already use the declared expected-utility scale. Subtracting a
constant cost below means a cost measured on that same scale.
\index{decision policy!randomized}\index{value of information}
\index{perfect information}
\end{hbdefinition}

\begin{hbtheorem}{thm:p17-information-policy-value}{proved=annex:p17.information-policy-value}{Optimal policies and the value formula}
An optimal deterministic policy exists and randomization cannot improve
the optimum. Specifically,
\[
 V(\pi,P,u)
 =\sum_{z=1}^r\max_{a\in A}\sum_s\pi_s p_{sz}u(a,s)
 =\sum_{z:q_z>0}q_z\max_{a\in A}\sum_s\pi^z_s u(a,s).
\]
At each positive-probability signal, select an action maximizing its
posterior expected payoff. At a zero-probability signal any policy may be
specified; it has no effect on the ex ante value.
\end{hbtheorem}

\begin{hbtheorem}{thm:p17-information-value-bounds}{proved=annex:p17.information-value-bounds}{Information value, perfect-information bound and cost}
For every finite problem,
\[
 V_0(\pi,u)\leq V(\pi,P,u)\leq V_{\mathrm{PI}}(\pi,u).
\]
Thus the expected value of sample information
\(\operatorname{EVSI}=V(\pi,P,u)-V_0(\pi,u)\) is nonnegative and at
most \(\operatorname{EVPI}=V_{\mathrm{PI}}(\pi,u)-V_0(\pi,u)\).
If using the experiment incurs a fixed cost \(c\geq0\) in payoff units,
its net value is \(V(\pi,P,u)-c\). Acquisition is weakly preferable to
the no-information option exactly when \(c\leq\operatorname{EVSI}\).
These statements assume that the signal may be ignored and the original
actions remain available.
\end{hbtheorem}

\section{State-independent processing of a signal}
\label{sec:p17-garbling}

\begin{hbdefinition}{def:p17-garbling}{A stochastic transformation and the source's sufficiency comparison}
Let \(Q=(q_{sy})\) be another experiment with the same states and
nonempty finite signal set \(Y=\{1,\ldots,t\}\).
A \emph{garbling} of \(P\) into \(Q\) is a row-stochastic
\(r\times t\) matrix \(M\) such that
\[
                         PM=Q.
\]
The processor observes \(z\) and selects \(y\) according to \(M_{zy}\).
Its rule is the same in every state.
Following the source, write \(P>Q\) when such a transformation exists.
Here \(>\) names this source comparison; it does not mean a strict order.
The identity matrix shows \(P>P\).
\index{garbling}\index{Blackwell comparison!stochastic transformation}
\end{hbdefinition}

\begin{hbtheorem}{thm:p17-garbling-simulation}{proved=annex:p17.garbling-simulation}{Simulation of every decision and composition}
If \(PM=Q\), any policy \(\delta(a\mid y)\) after \(Q\) is simulated
after \(P\) by
\[
                    \widehat\delta(a\mid z)
                      =\sum_y M_{zy}\delta(a\mid y).
\]
For every state separately, the action distribution and expected payoff
are identical. Consequently \(V(\pi,P,u)\geq V(\pi,Q,u)\) for every
prior and payoff problem.

The relation \(>\) is reflexive and transitive. Mutual garbling defines
an equivalence relation on finite experiments with the same state set;
equivalent experiments need not have the same signal labels or number of
signals. The identity experiment, whose signal is the state, garbles to
every experiment. Every experiment garbles to the single-signal experiment.
\end{hbtheorem}

\begin{hbdefinition}{def:p17-blackwell-risk-sets}{Loss vectors and the source's informativeness comparison}
Let \(\mathcal A\subseteq\mathbb R^n\) be nonempty, closed, bounded and
convex. A choice \(a_z\in\mathcal A\) at each signal gives the statewise
loss vector \(v\), with
\[
 v_s=\sum_z p_{sz}(a_z)_s.
\]
Let \(B(P,\mathcal A)\) be the set of all these vectors.
The convex set \(\mathcal A\) already permits randomization among its
points. In particular, a finite list of loss vectors gives its convex hull.
The empty set would give empty attainable sets for both experiments and
adds no restriction to the comparison.

Following the source, write \(P\supset Q\) when
\[
       B(P,\mathcal A)\supseteq B(Q,\mathcal A)
       \quad\hbox{for every such }\mathcal A.
\]
The symbol between experiments is a named non-strict comparison.
The symbol between attainable sets has its ordinary set-inclusion meaning
\cite[pp.~265--267]{Blackwell1953EquivalentExperiments}.
\index{Blackwell comparison!attainable loss vectors}
\end{hbdefinition}

\section{The finite comparison theorem}
\label{sec:p17-finite-blackwell}

\begin{hbdefinition}{def:p17-comparison-geometry}{Finite-dimensional geometry used in the proof}
For real arrays of the same finite shape, put
\[
 \langle H,J\rangle=\sum_{i,j}H_{ij}J_{ij},\qquad
 d^2(H,J)=\langle H-J,H-J\rangle.
\]
Vectors are included as one-column arrays. A convex set contains every
segment between two of its points. We use the usual Euclidean meaning of
closed and bounded.
For a square matrix \(H\), define
\(\operatorname{Trace}(H)=\sum_i H_{ii}\) and
\(\operatorname{diag}(H)=(H_{11},\ldots,H_{nn})\).
The trace notation follows the source's matrix proof.
For a real number \(t\), write \(\operatorname{abs}(t)=\max(t,-t)\).
In the finite game below, \(H_{ij}\) is the first player's payoff and
\(-H_{ij}\) the second player's payoff. Their independent randomized
row and column choices have expected payoff \(\alpha^{\mathsf T}H\beta\).
\end{hbdefinition}

\begin{hblemma}{thm:p17-finite-separation-minimax}{proved=annex:p17.finite-separation-minimax}{Separation and finite minimax}
Let \(C\) be a nonempty closed bounded convex subset of a finite-dimensional
real space, and \(b\notin C\). There is \(c_0\in C\) minimizing
\(d^2(b,c)\), and, with \(h=b-c_0\),
\[
 \langle h,c\rangle\leq\langle h,c_0\rangle
                      <\langle h,b\rangle\qquad(c\in C).
\]
For every finite real payoff matrix \(H\), with nonempty row and column
sets, the corresponding finite zero-sum game satisfies
\[
 \max_{\alpha\in\Delta(\text{rows})}\min_\beta \alpha^{\mathsf T}H\beta
 =
 \min_{\beta\in\Delta(\text{columns})}\max_\alpha
                  \alpha^{\mathsf T}H\beta.
\]
Both optima exist and give a saddle point. The annex proves the needed
compactness consequences, separation and minimax argument.
\end{hblemma}

\begin{hbtheorem}{thm:p17-finite-blackwell}{proved=annex:p17.finite-blackwell}{Equivalent finite experiment comparisons}
For finite experiments \(P,Q\) with the same nonempty state set, the
following are equivalent:
\begin{enumerate}
\item \(P>Q\): some row-stochastic \(M\) satisfies \(PM=Q\).
\item \(P\supset Q\): every loss vector attainable with \(Q\) in every
      closed bounded convex loss set is attainable with \(P\).
\item \(V(\pi,P,u)\geq V(\pi,Q,u)\) for every prior and every nonempty
      finite action set with real payoffs.
\item For one specified full-support prior \(\pi^0\), the same value
      inequality holds for every finite action/payoff problem.
\end{enumerate}
For the last condition \(\pi^0_s>0\) for every state. The prior is fixed
across all decision problems. This is the finite form of Blackwell's
comparison and the source's Sherman--Stein theorem
\cite[theorems~3--6, pp.~267--268]{Blackwell1953EquivalentExperiments}.
\end{hbtheorem}

\begin{hbtheorem}{thm:p17-blackwell-trace-route}{proved=annex:p17.blackwell-trace-route}{The source's trace condition and minimax construction}
If \(P\) is \(n\times r\), \(Q\) is \(n\times t\), and \(P\supset Q\),
then for every real \(t\times n\) matrix \(D\) there is a row-stochastic
\(r\times t\) matrix \(M\) for which
\[
       \operatorname{diag}(PMD)=\operatorname{diag}(QD).
\]
In particular their traces agree.
Moreover, the weaker condition
\[
 \forall D\in[0,1]^{t\times n}\ \exists M\text{ row stochastic}:
       \operatorname{Trace}((Q-PM)D)\leq0
\]
already implies that one row-stochastic \(M_0\) satisfies \(PM_0=Q\).
The matrix \(M\) in the quantified premise may depend on \(D\); the
conclusion supplies one matrix satisfying all the experiment equations.
The annex gives the full finite minimax argument behind source Theorem~6.
\end{hbtheorem}

\begin{hbmethod}{def:p17-garbling-certificate}{Checking a proposed information comparison}
For supplied finite matrices \(P,Q\), a proposed certificate \(M\) must
have exactly \(r\times t\) entries, all nonnegative, with every row sum
one and every entry of \(PM\) equal to the corresponding entry of \(Q\).
These are finite linear constraints. A valid certificate gives the explicit
simulation of every decision rule.

The separation proof supplies an alternative certificate when no such
matrix exists: a real array whose linear value at \(Q\) exceeds its value
at every garbling \(PM\). The proof converts that array into a decision
problem that strictly favors \(Q\). This characterizes the certificates;
it is not a claim that arbitrary unspecified real inputs admit an effective
finite-precision decision procedure.
\end{hbmethod}

\section{Posterior beliefs and mean-preserving comparison}
\label{sec:p17-posterior-comparison}

\begin{hbdefinition}{def:p17-posterior-law}{The finite distribution of posterior beliefs}
For experiment \(P\) and prior \(\pi\), its posterior law assigns mass
\(q_z\) to \(\pi^z\) for each \(q_z>0\), combining repeated posterior
vectors. Denote it by \(\mu_{P,\pi}\). Zero-probability signals contribute
no point to this law.
A real function \(\phi\) on a convex set is convex when
\[
 \phi(\lambda x+(1-\lambda)y)
       \leq\lambda\phi(x)+(1-\lambda)\phi(y)
       \quad(0\leq\lambda\leq1).
\]
All posterior expectations below are finite sums.
\index{posterior beliefs!distribution}\index{Bayes plausibility}
\end{hbdefinition}

\begin{hbtheorem}{thm:p17-posterior-splitting}{proved=annex:p17.posterior-splitting}{Posterior averaging and implementation}
The posterior law has barycenter \(\pi\):
\[
                         \sum_{z:q_z>0}q_z\pi^z=\pi.
\]
Conversely, let a finite list of probability vectors \(\nu^z\) have
weights \(q_z\geq0\) summing to one, and suppose
\(\sum_z q_z\nu^z=\pi\). Some finite experiment implements the given
positive-weight posterior law. On every state with \(\pi_s>0\), take
\[
                         p_{sz}=q_z\nu^z_s/\pi_s.
\]
Rows at states with zero prior probability may be any probability vector.
The barycenter equation guarantees that positive-weight posteriors assign
zero probability to such states.
\end{hbtheorem}

\begin{hbtheorem}{thm:p17-finite-posterior-order}{proved=annex:p17.finite-posterior-order}{Finite posterior convex order}
Fix a full-support prior \(\pi\). Then \(P>Q\) exactly when
\[
 \sum_z q_z\phi(\pi^z)
       \geq\sum_y \widetilde q_y\phi(\widetilde\pi^y)
\]
for every finite-valued continuous convex \(\phi\) on \(\Delta(S)\).
Only positive-weight signals are included.
In fact the converse needs only functions that are maxima of finitely
many linear functions.

If \(Q=PM\), the coupled signals have masses \(q_zM_{zy}\). At every
positive-probability output \(y\),
\[
       \widetilde\pi^y
          =\sum_z\frac{q_zM_{zy}}{\widetilde q_y}\,\pi^z.
\]
Thus processing averages the earlier posteriors conditional on the later
signal. This proves the direction of the convex-order comparison, rather
than treating it as a convention about which distribution is more spread out.
\end{hbtheorem}

\begin{hbtheorem}{thm:p17-finite-convex-transport}{proved=annex:p17.finite-convex-transport}{Finite convex order and a mean-preserving transformation}
Let \(X\) be a nonempty finite subset of \(\mathbb R^d\), \(d\geq1\),
and let \(\mu,\nu\) be probability masses on \(X\). Then the following
are equivalent:
\begin{enumerate}
\item \(\sum_x\mu_x\phi(x)\geq\sum_x\nu_x\phi(x)\) for every continuous
      convex function \(\phi\) on the convex hull of \(X\).
\item There is a row-stochastic matrix \(T=(T_{yx})\) on \(X\) such that
      \[
       \sum_y\nu_yT_{yx}=\mu_x,\qquad
       \sum_xT_{yx}x=y\quad\hbox{for every }y\in X.
      \]
\end{enumerate}
The second equality is required even in rows with \(\nu_y=0\).
The proof constructs those rows explicitly and supplies the affine
embedding left implicit in the source's finite-vector reduction
\cite[theorem~7, pp.~267--268]{Blackwell1953EquivalentExperiments}.
\end{hbtheorem}

\section{Calculations and limits of an information claim}
\label{sec:p17-information-examples}

\begin{hbexample}{def:p17-binary-signal-processing}{An explicit garbling certificate}
For two states and two signals, let
\[
 P=\begin{pmatrix}4/5&1/5\\1/5&4/5\end{pmatrix},
 \quad
 M=\begin{pmatrix}3/4&1/4\\1/4&3/4\end{pmatrix},
 \quad
 Q=\begin{pmatrix}13/20&7/20\\7/20&13/20\end{pmatrix}.
\]
Then \(PM=Q\). The processor independently flips the signal label with
probability \(1/4\). Under a uniform prior and unit payoff for a correct
state guess, the values are \(4/5\) and \(13/20\), respectively.
Thus the garbling comparison holds for every decision problem, while this
particular problem also demonstrates a strict loss and excludes a reverse
garbling.
\end{hbexample}

\begin{hbexample}{def:p17-service-information}{A signal that changes an action}
A device is in state bad with probability \(1/5\) and good with
probability \(4/5\). Operating has payoff \(-10\) in the bad state and
zero in the good state; servicing has payoff \(-3\) in either state.
An alert occurs with probability \(4/5\) in the bad state and \(1/10\)
in the good state.

Without the signal the best value is \(-2\), from operating. The alert
probability is \(6/25\), with bad-state posterior \(2/3\); the no-alert
posterior is \(1/19\). The optimal policy services after an alert and
operates otherwise. Its value is \(-28/25\), so
\[
               \operatorname{EVSI}=22/25,\qquad
               \operatorname{EVPI}=7/5.
\]
A signal cost \(1/2\) gives net value \(-81/50>-2\).
These probabilities and payoffs are stipulated mathematical data.
\end{hbexample}

\begin{hbexample}{def:p17-incomparable-experiments}{Two useful experiments that are not ordered}
On three equally likely states, consider
\[
 P=\begin{pmatrix}1&0\\0&1\\0&1\end{pmatrix},
 \qquad
 Q=\begin{pmatrix}0&1\\1&0\\0&1\end{pmatrix}.
\]
Experiment \(P\) reveals whether the state is 1; \(Q\) reveals whether
it is 2. With unit payoff for correctly answering whether the state is 1,
the respective values are \(1\) and \(2/3\). For the state-2 question,
the values reverse. Neither experiment garbles to the other.
A favorable value calculation for one task is therefore not a universal
information comparison.
\end{hbexample}

\begin{hbexample}{def:p17-finite-transport-example}{Splitting a mean into two outcomes}
Let \(X=\{-1,0,1\}\), let \(\mu\) give mass \(1/2\) to each endpoint,
and let \(\nu\) be concentrated at zero. The transformation sends zero
equally to the two endpoints and leaves each endpoint fixed in its
zero-weight input row. It sends \(\nu\) to \(\mu\) and preserves the mean
in every row. The convex-order inequality is exactly
\(\phi(0)\leq(\phi(-1)+\phi(1))/2\).
If instead \(\nu\) is concentrated at 1, the inequality already fails
for \(\phi(x)=x\). The required mean agreement is a consequence of convex
order, not a freely adjustable parameter of the transport.
\end{hbexample}

\begin{hbtheorem}{thm:p17-information-example-checks}{proved=annex:p17.information-example-checks}{The calculations and the zero-prior boundary}
The preceding examples have the stated probabilities, values and
comparison properties.
The full-support premise in the fixed-prior characterization cannot be
discarded: with two states, prior \((1,0)\), the single-signal experiment
and the identity experiment have equal value in every payoff problem at
that prior, but the single-signal experiment cannot garble to the identity.
More generally, fixed-prior value comparisons cannot identify experiment
rows belonging only to zero-prior states.
\end{hbtheorem}

\begin{proofboundary}{def:p17-information-next}{Further information and sequential-decision regimes}
All statements in this chapter have complete finite proofs.
\Cref{sec:p17-restricted-decisions-binary-testing} supplies the finite
restricted \(k\)-decision comparisons, binary testing-error formulation
and conditionally independent repetition theorem.
Blackwell's general bounded transport and the resulting observation-space
comparison are supplied in
\cref{thm:p17-general-convex-transport,thm:p17-general-decision-comparison,
thm:p17-general-garbling}, with their measurable-space hypotheses.
Sequential information choice requires a
specified history, feasible actions, signal laws and continuation payoff.
The static value formula alone does not establish an optimal adaptive policy,
dynamic consistency, or a value comparison robust to arbitrary additional
signals. Those are separate decision-theoretic treatments.
\end{proofboundary}
